% Vietnamese translation for OI Public Library
% Source-File: IOI中国国家候选队论文/2008/Day2/2.高逸涵《部分贪心思想在信息学竞赛中的应用》/部分贪心在信息学竞赛中的应用.doc
\documentclass[11pt]{article}
\usepackage{oipl-vn}

\newcommand{\Oh}{\mathcal{O}}

\title{Ứng dụng tư tưởng tham lam một phần trong thi tin học}
\author{Gao Yihan\\Trường Trung học trực thuộc Đại học Thanh Hoa}
\date{Trại đông Tin học toàn quốc, 2008}

\begin{document}
\maketitle

\origtitle{Ứng dụng tư tưởng tham lam một phần trong thi tin học}
\sourcepath{IOI China candidate team papers / 2008 / Day 2 / Gao Yihan / partial greedy.doc}

\renewcommand{\abstractname}{Tóm tắt}
\begin{abstract}
Trong một số bài toán có kích thước dữ liệu rất lớn, ta thường muốn dùng
thuật toán tham lam, nhưng tham lam thuần túy có thể gặp phản ví dụ. Khi đó
có thể dùng một phương án dung hòa: chỉ tham lam ở phần chắc chắn đúng để thu
nhỏ quy mô bài toán, rồi xử lý phần còn lại bằng tìm kiếm, quy hoạch động
hoặc liệt kê. Bài viết gọi cách làm này là tham lam một phần và minh họa qua
ba ví dụ.
\end{abstract}

\noindent\textbf{Từ khóa:} tham lam một phần.

\section{Dẫn nhập}

Tham lam là một tư tưởng rất quan trọng trong thi tin học. Nếu chứng minh
được tính đúng đắn, thuật toán tham lam thường có độ phức tạp tốt hơn rất
nhiều so với các cách làm tổng quát. Nó cũng thường có mã ngắn, ít bộ nhớ và
dễ cài đặt.

Khó khăn nằm ở chỗ: với phần lớn bài toán, ta khó tìm được một chiến lược tham
lam đơn giản và đúng hoàn toàn. Ngay cả khi một ý tưởng trông rất hợp lý, vẫn
cần chứng minh chặt chẽ; đôi khi chỉ vài trường hợp biên nhỏ cũng đủ làm nó
sai.

Tham lam một phần xuất hiện trong tình huống đó. Ta dùng tham lam ở đoạn giữa
của quá trình, nơi quy mô lớn và quy luật toàn cục đủ rõ; còn gần trạng thái
đầu hoặc trạng thái đích, nơi dễ phát sinh ngoại lệ nhỏ, ta dùng một thuật
toán chắc chắn đúng như liệt kê, tìm kiếm hoặc quy hoạch động. Cách làm này
cân bằng giữa hiệu suất của tham lam và độ chắc chắn của các thuật toán tổng
quát.

\section{Khái niệm cơ bản}

\subsection{Thuật toán tham lam}

Thuật toán tham lam đưa ra quyết định theo hướng tốt nhất ở thời điểm hiện
tại, rồi tiếp tục cho tới trạng thái cuối. Ví dụ trực quan trong bản gốc là
tìm đường từ \(S\) tới \(T\) trong một đồ thị nhiều lớp. Nếu mỗi bước chỉ chọn
cạnh ngắn nhất sang lớp kế tiếp, ta có thể đi theo \(S\to A\to E\to T\) với
tổng độ dài \(1+4+4=9\). Nhưng đường tốt hơn là
\(S\to B\to D\to T\), tổng độ dài \(2+2+2=6\). Lựa chọn cục bộ tốt nhất không
nhất thiết dẫn tới nghiệm toàn cục tốt nhất.

\subsection{Ưu và nhược điểm của tham lam}

Một thuật toán tham lam đúng thường rất hấp dẫn: ý tưởng trực tiếp, thời gian
nhanh, mã nguồn ngắn, bộ nhớ nhỏ. Nhưng nó có phạm vi áp dụng hẹp. Nếu không
có chứng minh, một chiến lược tham lam chỉ là trực giác; trong phòng thi, trực
giác ấy có thể sai bởi một cấu hình nhỏ khó nhìn thấy.

Các thuật toán tổng quát như quy hoạch động hoặc tìm kiếm thường dễ đảm bảo
đúng hơn, nhưng lại không chịu nổi dữ liệu quá lớn. Tham lam một phần đặt mục
tiêu nằm giữa hai cực này: dùng tham lam để giảm kích thước, rồi dùng thuật
toán chắc chắn đúng cho phần còn lại.

\subsection{Tham lam một phần}

Nhiều phản ví dụ của tham lam thuần túy có quy mô nhỏ. Ý tưởng tham lam được
thiết kế cho bức tranh lớn, còn các chi tiết nhỏ ở biên lại làm nó sai. Khi
đó ta có thể chọn một ngưỡng an toàn: nếu quy mô còn lớn hơn ngưỡng, thực hiện
bước tham lam đã chứng minh được; khi quy mô giảm xuống dưới ngưỡng, chuyển
sang liệt kê hoặc quy hoạch động.

Một hình ảnh đơn giản là tìm cực tiểu của một hàm mà thành phần chính phụ
thuộc mạnh vào \(x\), còn thành phần nhiễu chỉ ảnh hưởng trong phạm vi nhỏ.
Nếu chỉ lấy điểm cực tiểu của phần chính thì có thể sai, nhưng điểm cực tiểu
thật chắc chắn nằm gần đó. Ta tham lam chọn vùng gần cực tiểu của phần chính,
rồi liệt kê trong một lân cận nhỏ để lấy đáp án chính xác.

Tham lam một phần vì vậy là một cách mở rộng của tham lam. Nó không làm tăng
đáng kể độ phức tạp tổng thể, nhưng chuyển các trường hợp biên sang một cơ
chế dễ chứng minh hơn.

\section{Ví dụ 1: Camels}

\subsection{Mô tả}

Có \(p\) người, \(x\) gói nhỏ và \(y\) gói lớn cần đi qua sa mạc. Không ai
muốn đi bộ, nên cần dùng lạc đà chở người và hành lý. Mỗi con lạc đà chỉ được
chở một trong bốn kiểu sau:
\begin{itemize}
  \item nhiều nhất \(3\) gói nhỏ;
  \item nhiều nhất \(2\) gói lớn;
  \item \(1\) người và nhiều nhất \(2\) gói nhỏ;
  \item \(1\) người và \(1\) gói lớn.
\end{itemize}
Cần số lạc đà ít nhất. Giới hạn:
\[
  1\le p\le 10^9,\qquad 0\le x,y\le 10^9.
\]

Khi đã biết bao nhiêu người mang gói nhỏ và bao nhiêu người mang gói lớn, số
lạc đà còn lại có thể tính trực tiếp. Nếu liệt kê số người mang gói nhỏ thì
dữ liệu quá lớn; nếu dùng công thức đóng thì các phép làm tròn khiến phân
tích rắc rối.

\subsection{Tham lam an toàn}

Ta nhìn người như một phần của loại gói mà họ đang mang. Nếu một người mang
gói nhỏ, một lạc đà chở được người đó cùng \(2\) gói nhỏ, trong khi lạc đà
không chở người chở được \(3\) gói nhỏ. Nếu một người mang gói lớn, một lạc đà
chở được người đó cùng \(1\) gói lớn, trong khi lạc đà không chở người chở
được \(2\) gói lớn. Trực giác là: thường nên ưu tiên cho người mang gói nhỏ.

Bản gốc chứng minh một bước đủ an toàn: nếu còn hơn \(6\) gói nhỏ và còn ít
nhất \(2\) người, thì để người mang gói nhỏ không làm nghiệm xấu hơn. Nếu một
phương án không làm vậy mà cho mọi người còn lại mang gói lớn, ta có thể lấy
một lạc đà đang chở \(3\) gói nhỏ và hai lạc đà đang chở người với gói lớn,
đổi thành hai lạc đà chở người với gói nhỏ và một lạc đà chỉ chở gói lớn.
Phương án mới không dùng nhiều lạc đà hơn, lại còn chở được thêm một gói nhỏ.

Từ đó có thuật toán:
\begin{enumerate}
  \item nếu số gói nhỏ và số người đều còn lớn hơn một ngưỡng an toàn, thực
  hiện một loạt bước cho người mang gói nhỏ;
  \item khi số người hoặc số gói nhỏ đã nhỏ, liệt kê số người mang gói nhỏ và
  tính trực tiếp phần còn lại.
\end{enumerate}

Trong chương trình gốc, các ngưỡng được lấy dư dả hơn điều kiện chứng minh,
chẳng hạn giữ lại vài người và vài chục gói nhỏ trước khi liệt kê. Đây là một
thủ pháp thường gặp: tăng ngưỡng một chút để tránh lỗi biên, trong khi độ phức
tạp vẫn không đổi.

\subsection{Nhận xét}

Bài này có các đặc điểm rất hợp với tham lam một phần: ràng buộc đơn giản,
miền dữ liệu rất lớn, ý tưởng tham lam rõ ràng, nhưng phản ví dụ có thể nằm ở
kích thước nhỏ. Tham lam làm nhiệm vụ giảm dữ liệu, còn liệt kê nhỏ bảo đảm
đúng.

\section{Ví dụ 2: Huge Knapsack}

\subsection{Mô tả}

Có \(N\) loại tượng vàng, mỗi loại có vô hạn bản sao, thể tích \(V_i\) và
trọng lượng \(W_i\). Có \(S\) túi, mỗi túi dung tích \(Y\). Mỗi lần có thể
trả chi phí \(C\) để hợp hai túi thành một túi lớn hơn; hợp \(3\) túi thành
một tốn \(2C\). Cần tối đa hóa tổng lượng vàng lấy được trừ chi phí hợp túi.
Giới hạn trong đề rất lớn:
\[
  S,Y\le 10^9,\qquad N\le 1000,\qquad W_i\le 100,\qquad V_i\le 18.
\]

Trước hết, với cùng thể tích chỉ cần giữ loại có trọng lượng lớn nhất, nên số
loại còn lại không vượt quá \(18\). Đây là bài ba lô vô hạn, nhưng dung tích
túi quá lớn để dùng quy hoạch động trực tiếp.

\subsection{Giảm dung tích túi}

Chọn loại có tỉ lệ \(W_i/V_i\) tốt nhất làm vật phẩm chủ đạo. Khi dung tích
còn lại đủ lớn, một nghiệm tối ưu phải tiếp tục dùng vật phẩm chủ đạo; nếu một
loại khác xuất hiện quá nhiều, ta có thể thay một cụm của nó bằng vật phẩm
chủ đạo mà không kém hơn.

Bản gốc dùng một chặn \(Limit\) nhỏ: với mỗi loại không chủ đạo, số lượng của
nó trong phần dư không cần vượt quá thể tích của vật phẩm chủ đạo. Vì
\(V_i\le 18\), có thể lấy \(Limit < 324\). Từ đó, quy hoạch động chỉ cần xử lý
phần dung tích dưới \(Limit\). Với dung tích \(k\) lớn, đáp án bằng:
\[
  \text{nhiều vật phẩm chủ đạo} + \text{giá trị tốt nhất của phần dư nhỏ}.
\]
Nhờ vậy truy vấn giá trị tốt nhất cho một túi có thể làm trong thời gian hằng
số sau tiền xử lý.

\subsection{Giảm số túi cần hợp}

Vẫn còn tham số \(S\) rất lớn. Ta tiếp tục dùng cùng tư tưởng. Vì phần dư theo
modulo \(Limit\) chỉ có \(Limit\) khả năng, không cần hợp quá \(Limit\) túi
trong một cụm. Nếu hợp \(A\) túi và \(B\) túi tạo ra cùng lớp dư với
\(A>B\), phần chênh \(A-B\) có thể được lấp đầy bằng vật phẩm chủ đạo mà không
cần tham gia cụm hợp túi. Do đó việc hợp quá nhiều túi không đem lại lợi ích
thêm.

Bài toán còn lại giống một bài ba lô có dung lượng \(S\) nhưng chỉ có một số
lượng nhỏ kiểu cụm cần xét, và có thể giải bằng cùng kỹ thuật trên. Ở đây
tham lam một phần được dùng hai lần: một lần để giảm \(Y\), một lần để giảm
\(S\).

\subsection{Nhận xét}

Tham lam thuần túy ``luôn lấy vật phẩm có tỉ lệ tốt nhất'' không đúng ở phần
dư nhỏ. Quy hoạch động thuần túy thì không chịu nổi dung tích lớn. Phối hợp
hai cách giúp bài toán giữ được độ đúng của quy hoạch động ở miền nhỏ và tốc
độ của tham lam ở miền lớn.

\section{Ví dụ 3: Cow Relays}

\subsection{Mô tả}

Cho một đồ thị vô hướng có \(T\) cạnh, mỗi đỉnh kề ít nhất hai cạnh, mỗi cạnh
có độ dài. Hỏi đường đi ngắn nhất từ đỉnh xuất phát tới đỉnh kết thúc với đúng
\(N\) cạnh là bao nhiêu. Giới hạn:
\[
  1\le N\le 10^9,\qquad T\le 100.
\]
Đồ thị rất nhỏ nhưng số cạnh cần đi rất lớn, đây là tín hiệu điển hình để tìm
một phần có thể xử lý bằng tham lam hoặc bằng một cấu trúc lặp.

\subsection{Cấu trúc của đường đi tối ưu}

Trực giác là: để đường đi có đúng \(N\) cạnh mà tổng độ dài nhỏ nhất, ta sẽ đi
qua lại rất nhiều lần trên một cạnh ngắn. Trước hết cần chứng minh rằng phần
``đi qua lại nhiều lần'' có thể tập trung trên một cạnh, không cần phân tán ở
nhiều cạnh.

Xét một đường đi tối ưu và chọn cạnh ngắn nhất xuất hiện trên nó. Nếu một cạnh
khác bị đi qua lại nhiều lần, ta xóa một cặp đi-về trên cạnh đó và thay bằng
một cặp đi-về trên cạnh ngắn nhất. Số cạnh không đổi, điểm đầu và điểm cuối
không đổi, tổng độ dài không tăng. Vì vậy có thể giả sử mọi lần lặp lớn đều
nằm trên cạnh ngắn nhất được chọn.

Tiếp theo, bản gốc chứng minh rằng ngoài cạnh dùng để lặp, không cần đi liên
tiếp cùng hướng qua một cạnh ba lần. Nếu điều này xảy ra, ta có thể tách được
một hoặc hai chu trình từ đoạn đường đã đi; khi chu trình có độ dài cạnh chẵn,
xóa chu trình đó và thay bằng các lần đi-về trên cạnh ngắn nhất. Nếu hai chu
trình đều lẻ, xóa cả hai để thu được số cạnh chẵn và thay tương tự. Tổng độ
dài không tăng.

Kết quả là: ngoài cạnh ngắn nhất được dùng để lặp, mỗi cạnh khác chỉ cần xuất
hiện một số lần có chặn, cụ thể không quá bậc hằng theo \(T\). Phần không
tham lam của đường đi vì vậy chỉ có độ dài \(\Oh(T)\).

\subsection{Thuật toán}

Ta dùng quy hoạch động để tính đường ngắn nhất từ điểm đầu và từ điểm cuối
tới mọi đỉnh với số cạnh \(k\), trong đó \(0\le k\le 2T\). Sau đó liệt kê
cạnh được chọn làm cạnh lặp chính và chia số cạnh còn lại cho đoạn đầu, đoạn
lặp, đoạn cuối. Phần lặp có thể tính trực tiếp bằng độ dài cạnh ấy nhân với
số lần đi-về cần thiết, còn hai đầu đường được lấy từ bảng quy hoạch động.

Lời giải chính thức của bài dùng nhân ma trận/bình phương nhanh trên đại số
min-plus. Bản gốc nhấn mạnh rằng tham lam một phần cho một cách nhìn khác:
nó nén tham số khổng lồ \(N\) xuống một miền nhỏ quanh \(2T\), rồi để quy
hoạch động xử lý phần còn lại.

\section{Tổng kết}

Ba ví dụ thể hiện ba dạng dùng tham lam một phần. Trong \textit{Camels}, nó
hòa giải giữa một chiến lược tham lam hiệu quả và nhiều ngoại lệ nhỏ. Trong
\textit{Huge Knapsack}, nó giảm hai tham số rất lớn về miền có thể quy hoạch
động. Trong \textit{Cow Relays}, nó phát hiện phần lặp lớn của đường đi tối ưu
và chỉ để lại một phần ngắn cần xử lý chính xác.

Từ bản chất, tham lam một phần là sự kết hợp giữa tham lam và một thuật toán
khác. Nó giữ được trực giác, mã ngắn và hiệu suất cao của tham lam, đồng thời
giải quyết được một số bài mà tham lam thuần túy không đủ chắc chắn. Dù không
phải một kỹ thuật ``thuần'' theo nghĩa sách giáo khoa, nó rất thực dụng trong
thi tin học: khi thuật toán chính thống quá khó hoặc quá nặng, tham lam một
phần có thể là một lựa chọn đáng cân nhắc.

\section*{Tài liệu tham khảo}

Liu Rujia và Huang Liang, \emph{Nghệ thuật thuật toán và thi tin học}, Nhà
xuất bản Đại học Thanh Hoa.

\section*{Ghi chú về phụ lục nguồn}

DOC gốc kèm chương trình cho ví dụ \textit{Cow Relays}. Theo quy ước của bản
dịch này, phần mã mẫu dài không được chuyển nguyên văn; các ý tưởng cần thiết
đã được diễn giải trong phần thuật toán.

\end{document}
